\section*{Chapter \ref{ch:iv:concurrency-operating-systems-and-networks} --- Concurrency, Operating Systems and Networks}

\begin{solution}{iq:iv:concurrency-operating-systems-and-networks:1}
Each \texttt{++counter} is a load, an add and a store; when the threads interleave between the load and the store, one increment
overwrites the other and is lost. The standard says more: unsynchronised concurrent access with a write is a data race and the
behaviour is undefined, so no output can be predicted, and ThreadSanitizer reports ``data race'' on this program (the chapter's
test asserts the diagnosis, never the number printed). Fix with an atomic, a mutex, or per-thread counters.
\iqlookfor{the read-modify-write mechanism and the formal status of a data race.}
\end{solution}

\begin{solution}{iq:iv:concurrency-operating-systems-and-networks:2}
When the critical section is a few instructions, the lock is rarely contended, and the waiting thread runs on its own core with
nothing better to do, as on isolated cores of a trading engine: spinning avoids a sleep and a wake-up through the kernel. Never
when the holder can be descheduled (a spinner then burns its whole time slice) or when many threads contend.
\iqlookfor{the trade-off between spinning and sleeping, tied to core isolation.}
\end{solution}

\begin{solution}{iq:iv:concurrency-operating-systems-and-networks:3}
The kernel saves one thread's registers and loads another's. Beyond that, the incoming thread finds caches, branch predictors and
translation buffers holding the other thread's state, so it runs slowly for a while. Trading systems pin hot threads to isolated
cores, keep them busy polling rather than blocking, and move non-critical work to other cores (One Quant Book~13, chapter~13).
\iqlookfor{direct and indirect costs, and pinning and isolation as the remedy.}
\end{solution}

\begin{solution}{iq:iv:concurrency-operating-systems-and-networks:4}
Market data goes to many receivers at once, and a late update is worthless: multicast UDP sends once to all subscribers and never
blocks on a lost packet; receivers detect gaps by sequence number and recover separately. Orders must arrive, once and in order, and
a loss must be known: TCP gives a reliable ordered stream and a session whose failure is visible.
\iqlookfor{fan-out and staleness against reliability and order.}
\end{solution}

\begin{solution}{iq:iv:concurrency-operating-systems-and-networks:5}
Two (for a million increments each). Thread A loads 0 and stalls; thread B completes all but its last increment; A stores 1,
overwriting B's work; B loads 1 for its last increment and stalls; A completes its remaining increments; B stores 2. An exhaustive
enumeration of all interleavings for small counts (\cref{prop:iv:concurrency-operating-systems-and-networks:counter}) confirms that
2 is the minimum. In C++ the program is a data race, so even this analysis describes the machine model, not a guarantee of the
language.
\iqlookfor{the adversarial interleaving and the caveat about undefined behaviour.}
\end{solution}

\begin{solution}{iq:iv:concurrency-operating-systems-and-networks:6}
Under sequential consistency no: some write comes first in the single order, so the other thread's later read sees it; the chapter's
explorer finds only (0,1), (1,0) and (1,1). On x86, yes: each store can sit in its core's store buffer while the load after it reads
memory. With C++ relaxed (or even acquire and release) atomics, yes. To forbid it, make all four operations \texttt{seq\_cst} or put a
\texttt{seq\_cst} fence between each store and the following load (on x86, an \texttt{mfence} or a locked instruction).
\iqlookfor{store buffering, the difference between SC and TSO, and the full fence.}
\end{solution}

\begin{solution}{iq:iv:concurrency-operating-systems-and-networks:7}
The flag store is a \texttt{release} and the flag load an \texttt{acquire}: every write before the release (the message) is then
visible to a thread whose acquire reads the flag's new value. It forbids ``flag seen, stale message'', the outcome (1, 0) of
the chapter's message-passing test, which sequential consistency also forbids and relaxed ordering allows. This is the pattern of
the SPSC queue's indices.
\iqlookfor{the release--acquire pairing and the outcome it excludes.}
\end{solution}

\begin{solution}{iq:iv:concurrency-operating-systems-and-networks:8}
Eight 8-byte counters fill one 64-byte cache line, so each write invalidates the line in the other core's cache and the line
ping-pongs between cores (false sharing, One Quant Book~13, chapter~3). Give each counter its own line (\texttt{alignas(64)} or padding),
or keep counters thread-local and combine them at the end.
\iqlookfor{cache-line granularity of coherence and padding as the fix.}
\end{solution}

\begin{solution}{iq:iv:concurrency-operating-systems-and-networks:9}
A page holds $4096/8 = 512$ entries, so each level resolves 9 bits; the offset takes 12 bits; $(48 - 12)/9 = 4$ levels. TLB reach:
$1536 \times 4\,\text{KiB} = 6\,\text{MiB}$ with small pages, $1536 \times 2\,\text{MiB} = 3\,\text{GiB}$ with huge pages, so a
working set of a few gigabytes stays translated with huge pages and misses constantly without them.
\iqlookfor{the radix arithmetic and the reach argument for huge pages.}
\end{solution}

\begin{solution}{iq:iv:concurrency-operating-systems-and-networks:10}
\Cref{lst:iv:concurrency-operating-systems-and-networks:spsc}. The producer reads its own \texttt{tail} relaxed (only it writes it),
reads \texttt{head} with acquire (so the consumer's reads of a slot happen before the producer overwrites it), writes the element,
then stores \texttt{tail} with release (publishing the element). The consumer mirrors it: acquire on \texttt{tail}, read, release on
\texttt{head}. The indices sit on separate cache lines to avoid false sharing. The chapter's stress test passes two million sequence
numbers in order and runs clean under ThreadSanitizer.
\iqlookfor{each ordering justified by what it publishes, and the cache-line layout.}
\end{solution}

\begin{solution}{iq:iv:concurrency-operating-systems-and-networks:11}
The card receives the frame, checks it and writes it by DMA into a ring of buffers in host memory; it raises an interrupt, the
driver schedules polling (NAPI) and hands the buffer to the stack; IP and UDP headers are processed, the socket is looked up and the
payload queued on its buffer; a thread blocked in \texttt{recvmsg} is woken, and the system call copies the payload to user memory
(\cref{fig:iv:concurrency-operating-systems-and-networks:packet}). Kernel bypass maps the card's ring into the process, which polls it:
the interrupt, the stack traversal, the wake-up, the system call and the copy all disappear.
\iqlookfor{the full path, and exactly which steps bypass removes.}
\end{solution}

\begin{solution}{iq:iv:concurrency-operating-systems-and-networks:12}
\omcode{code/interviews/25-concurrency-operating-systems-and-networks/python/iv_conc.py}{76}{102}{A gap detector that tolerates
duplicates and fills gaps with late arrivals.}\label{lst:iv:concurrency-operating-systems-and-networks:gaps}
Track the next expected number; a higher one opens a gap, a lower one is a duplicate or a late arrival that closes part of a gap.
When a gap stays open past a timeout, take the packet from the redundant line (arbitration), request a retransmission, or, for a
book feed, mark the book invalid and rebuild from a snapshot (One Quant Book~10, chapter~26). The chapter's test compares the
detector with a direct count of missing numbers on 300 random streams.
\iqlookfor{gap bookkeeping with late packets, and the recovery hierarchy.}
\end{solution}

\begin{solution}{iq:iv:concurrency-operating-systems-and-networks:13}
Deadlock: A holds the book and waits for the risk state while B holds the risk state and waits for the book. Prevent it by a global
lock order (always the book first), or by acquiring both at once with a deadlock-avoiding primitive (\texttt{std::scoped\_lock} on both
mutexes); better still, design so that one thread owns each piece of state and others send it messages.
\iqlookfor{the circular wait and the standard preventions.}
\end{solution}
